% !TEX program = xelatex
\documentclass[11pt,a4paper]{article}
\usepackage[margin=2.2cm,top=2cm]{geometry}
\usepackage{fontspec}
\usepackage{polyglossia}
\setdefaultlanguage{turkish}
\usepackage{amsmath,amssymb}
\usepackage{graphicx}
\usepackage{booktabs}
\usepackage{enumitem}
\usepackage{multicol}
\usepackage[most]{tcolorbox}
\usepackage{fancyhdr}
\usepackage{titlesec}
\usepackage{xcolor}
\usepackage{needspace}
\usepackage[hidelinks]{hyperref}

\definecolor{anamavi}{RGB}{0,45,98}
\definecolor{cozumyesil}{RGB}{0,110,50}
\definecolor{uyarikirmizi}{RGB}{180,20,20}

\titleformat{\section}{\large\bfseries\color{anamavi}}{\thesection.}{0.5em}{}
\titleformat{\subsection}{\normalsize\bfseries\color{anamavi}}{\thesubsection}{0.5em}{}

\pagestyle{fancy}\fancyhf{}
\lhead{\small MMAT105 Genel Matematik}
\rhead{\small Hafta 11 -- Ders Notu}
\cfoot{\small\thepage}
\renewcommand{\headrulewidth}{0.4pt}

\newtcolorbox{kural}[1][Kural]{colback=anamavi!5,colframe=anamavi,title=#1,fonttitle=\bfseries,breakable,left=4pt,right=4pt,top=2pt,bottom=2pt}
\newtcolorbox{dikkat}[1][Dikkat]{colback=uyarikirmizi!5,colframe=uyarikirmizi,title=#1,fonttitle=\bfseries,breakable,left=4pt,right=4pt,top=2pt,bottom=2pt}
\newcounter{ornek}
\newenvironment{ornek}{\refstepcounter{ornek}\par\medskip\noindent\textbf{\color{anamavi}Örnek \theornek.}\ }{\par}
\newenvironment{cozum}{\begin{tcolorbox}[colback=cozumyesil!6,colframe=cozumyesil,boxrule=0.5pt,left=4pt,right=4pt,top=2pt,bottom=2pt,breakable]\textbf{\color{cozumyesil}Çözüm.}\ }{\end{tcolorbox}}
\newcounter{alis}
\newenvironment{alistirma}{\refstepcounter{alis}\par\medskip\needspace{9\baselineskip}\noindent\textbf{\color{anamavi}Alıştırma \thealis.}\ }{\par}
\newcommand{\kendinizoncesi}{\needspace{12\baselineskip}}
\newcommand{\ab}[1]{\overline{#1}}
\setlength{\parindent}{0pt}
\setlength{\parskip}{3pt}

\begin{document}


% ---------------- Başlık ----------------
\begin{minipage}[c]{\textwidth}
{\LARGE\bfseries\color{anamavi} MMAT105 Genel Matematik}\\[4pt]
{\large Hafta 11 Ders Notu: Çarpanlara Ayırma ve Özdeşlikler\\ (İki Kare Farkı, Tam Kare, Küp Açılımları)}\\[6pt]
{\small 2026--2027 Güz}
\end{minipage}

\vspace{4pt}\noindent\textcolor{anamavi}{\rule{\textwidth}{1pt}}

\textbf{Bu notun kullanımı.} Her bölümde önce kısa bir kural kutusu, ardından çözümlü örnekler yer alır. Örnekleri çözümüne bakmadan önce kendiniz deneyin. Notun sonunda önce çözümlü alıştırmalar, ardından cevap anahtarlı ``Kendiniz Çözün'' alıştırmaları yer alır.

% =====================================================
\section{Ortak Çarpan Parantezine Alma ve Gruplandırma}

\begin{kural}[Kural]
Bir ifadeyi \textbf{çarpanlarına ayırmak}, onu çarpım biçiminde yazmaktır.
\begin{itemize}[leftmargin=*,itemsep=1pt]
\item Ortak çarpan: \(A\cdot x+A\cdot y-A\cdot z=A\,(x+y-z)\). Ortak çarpan, katsayıların EBOB'u ile her terimde bulunan harflerin en küçük kuvvetlerinin çarpımıdır.
\item Gruplandırma: \(ax+ay+bx+by=a(x+y)+b(x+y)=(x+y)(a+b)\).
\end{itemize}
\end{kural}

\begin{dikkat}
Sonucu kontrol etmek için parantezleri geri dağıtın. \(b-a=-(a-b)\). Bir grubun önündeki \(-\) işareti parantez dışına alınınca içerideki işaretler değişir: \(-3x-6=-3(x+2)\).
\end{dikkat}

\begin{ornek}
\(6x^3y-9x^2y^2+3xy\) ifadesini çarpanlarına ayırınız.
\end{ornek}
\begin{cozum}
EBOB\((6,9,3)=3\); her terimde en az \(x^1\) ve \(y^1\) var \(\Rightarrow\) ortak çarpan \(3xy\).
\[6x^3y-9x^2y^2+3xy=3xy\,(2x^2-3xy+1)\]
Son terimden geriye \(1\) kalır; bu \(1\) yazılmazsa sonuç yanlış olur.
\end{cozum}

\begin{ornek}
\(x^3-2x^2+5x-10\) ifadesini çarpanlarına ayırınız.
\end{ornek}
\begin{cozum}
\[x^3-2x^2+5x-10=x^2(x-2)+5(x-2)=(x-2)(x^2+5)\]
\end{cozum}

\begin{ornek}
\(x+y=5\) ve \(x^2y+xy^2=30\) olduğuna göre \(x\cdot y\) kaçtır?
\end{ornek}
\begin{cozum}
\(x^2y+xy^2=xy\,(x+y)\) \(\Rightarrow\) \(30=xy\cdot 5\) \(\Rightarrow\) \(xy=6\).
\end{cozum}

% =====================================================
\section{İki Kare Farkı}

\begin{kural}[Özdeşlik]
\[a^2-b^2=(a-b)(a+b)\]
Üsler çift ise ardışık uygulanabilir: \(x^4-y^4=(x^2-y^2)(x^2+y^2)=(x-y)(x+y)(x^2+y^2)\).
\end{kural}

\begin{dikkat}
\(a^2+b^2\) (iki kare toplamı) reel sayılarda bu yolla çarpanlarına ayrılmaz. \((a-b)^2\) ile \(a^2-b^2\) farklıdır.
\end{dikkat}

\begin{ornek}
\((2x+1)^2-(x-2)^2\) ifadesini çarpanlarına ayırınız.
\end{ornek}
\begin{cozum}
\(a=2x+1\), \(b=x-2\) alınır:
\[\big[(2x+1)-(x-2)\big]\big[(2x+1)+(x-2)\big]=(x+3)(3x-1)\]
\end{cozum}

\begin{ornek}
\(x^8-1\) ifadesini çarpanlarına ayırınız.
\end{ornek}
\begin{cozum}
\[x^8-1=(x^4-1)(x^4+1)=(x^2-1)(x^2+1)(x^4+1)=(x-1)(x+1)(x^2+1)(x^4+1)\]
\end{cozum}

\begin{ornek}
\(37\cdot 43\) \ ve \ \(2026\cdot 2024-2025^2\) işlemlerinin sonucu kaçtır?
\end{ornek}
\begin{cozum}
\(37\cdot 43=(40-3)(40+3)=1600-9=1591\).\\
\(2026\cdot 2024=(2025+1)(2025-1)=2025^2-1\) \(\Rightarrow\) \(2026\cdot 2024-2025^2=-1\).
\end{cozum}

\begin{ornek}
\(\dfrac{6}{\sqrt{5}-\sqrt{2}}\) ifadesinin paydasını rasyonel yapınız.
\end{ornek}
\begin{cozum}
Pay ve payda eşleniği \(\sqrt{5}+\sqrt{2}\) ile çarpılır; paydada iki kare farkı oluşur:
\[\frac{6(\sqrt{5}+\sqrt{2})}{(\sqrt{5})^2-(\sqrt{2})^2}=\frac{6(\sqrt{5}+\sqrt{2})}{3}=2\sqrt{5}+2\sqrt{2}\]
\end{cozum}

\begin{ornek}
\(x^2-y^2=40\) ve \(x+y=10\) olduğuna göre \(x\) ve \(y\) kaçtır?
\end{ornek}
\begin{cozum}
\(40=(x-y)\cdot 10\) \(\Rightarrow\) \(x-y=4\). \ \(x+y=10\) ile toplanırsa \(2x=14\): \ \(x=7\), \(y=3\).
\end{cozum}

% =====================================================
\section{Tam Kare Özdeşlikleri}

\begin{kural}[Özdeşlikler]
\[(a+b)^2=a^2+2ab+b^2,\qquad (a-b)^2=a^2-2ab+b^2\]
\[a^2+b^2=(a+b)^2-2ab=(a-b)^2+2ab,\qquad (a+b)^2-(a-b)^2=4ab\]
\[x\neq 0\ \text{için}\quad x^2+\frac{1}{x^2}=\Big(x+\frac{1}{x}\Big)^2-2=\Big(x-\frac{1}{x}\Big)^2+2\]
\(x^2+bx+m\) ifadesinin tam kare olması için \(m=\big(\tfrac{b}{2}\big)^2\) olmalıdır.
\end{kural}

\begin{dikkat}
\((a+b)^2\neq a^2+b^2\). \ Ayrıca \(A^2=9\) ise \(A=3\) veya \(A=-3\)'tür; işaret bilgisi verilmedikçe tek cevap yoktur.
\end{dikkat}

\begin{ornek}
\(998^2\) sayısını tam kare özdeşliği ile hesaplayınız.
\end{ornek}
\begin{cozum}
\(998^2=(1000-2)^2=1000000-4000+4=996004\).
\end{cozum}

\begin{ornek}
\(a+b=5\) ve \(a^2+b^2=13\) olduğuna göre \(a\cdot b\) kaçtır?
\end{ornek}
\begin{cozum}
\((a+b)^2=a^2+b^2+2ab\) \(\Rightarrow\) \(25=13+2ab\) \(\Rightarrow\) \(ab=6\).
\end{cozum}

\begin{ornek}
\(x>0\) ve \(x^2+\dfrac{1}{x^2}=7\) olduğuna göre \(x+\dfrac{1}{x}\) kaçtır?
\end{ornek}
\begin{cozum}
\(\Big(x+\dfrac{1}{x}\Big)^2=x^2+\dfrac{1}{x^2}+2=9\). \ \(x>0\) olduğundan \(x+\dfrac{1}{x}>0\) \(\Rightarrow\) \(x+\dfrac{1}{x}=3\).
\end{cozum}

\begin{ornek}
\(9x^2+mx+16\) bir tam kare ve \(m<0\) ise \(m\) kaçtır?
\end{ornek}
\begin{cozum}
\(9x^2=(3x)^2\), \(16=4^2\). \(m<0\) olduğundan ifade \((3x-4)^2=9x^2-24x+16\) olmalıdır \(\Rightarrow\) \(m=-24\).
\end{cozum}

\begin{ornek}
\(x\) ve \(y\) reel sayılar ve \(x^2+y^2-4x+6y+13=0\) olduğuna göre \(x+y\) kaçtır?
\end{ornek}
\begin{cozum}
\(13=4+9\) olarak ayrılıp tam kareler oluşturulur:
\[(x^2-4x+4)+(y^2+6y+9)=0\ \Rightarrow\ (x-2)^2+(y+3)^2=0\]
Kareler negatif olamaz; toplamları \(0\) ise ikisi de \(0\)'dır: \(x=2\), \(y=-3\), \ \(x+y=-1\).
\end{cozum}

% =====================================================
\section{Üç Terimli İfadeler ve Çarpım Sıfır Kuralı}

\begin{kural}[Kural]
\begin{itemize}[leftmargin=*,itemsep=1pt]
\item \(x^2+bx+c=(x+m)(x+n)\), \quad burada \(m+n=b\) ve \(m\cdot n=c\).
\item \(ax^2+bx+c\): \(ax^2\) ve \(c\) çarpanlarına ayrılıp alt alta yazılır; çapraz çarpımların toplamı \(bx\) olacak biçimde seçilir.
\item \(A\cdot B=0\) ise \(A=0\) veya \(B=0\)'dır.
\end{itemize}
\end{kural}

\begin{ornek}
\(x^2-11x+28\) \ ve \ \(x^2+x-30\) ifadelerini çarpanlarına ayırınız.
\end{ornek}
\begin{cozum}
Çarpımı \(28\), toplamı \(-11\): \(-4\) ve \(-7\) \(\Rightarrow\) \(x^2-11x+28=(x-4)(x-7)\).\\
Çarpımı \(-30\), toplamı \(1\): \(6\) ve \(-5\) \(\Rightarrow\) \(x^2+x-30=(x+6)(x-5)\).
\end{cozum}

\begin{ornek}
\(3x^2-10x+8\) ifadesini çarpanlarına ayırınız.
\end{ornek}
\begin{cozum}
\begin{minipage}[t]{0.2\linewidth}
\begin{tabular}{ccc}
\(3x\) & & \(-4\)\\
\(x\) & & \(-2\)
\end{tabular}
\end{minipage}\hfill
\begin{minipage}[t]{0.62\linewidth}
Çapraz: \(3x\cdot(-2)+x\cdot(-4)=-10x\) \checkmark
\[3x^2-10x+8=(3x-4)(x-2)\]
\end{minipage}
\end{cozum}

\needspace{6\baselineskip}
\begin{ornek}
\(x^4-5x^2+4\) ifadesini çarpanlarına ayırınız.
\end{ornek}
\begin{cozum}
\(x^2\) tek bir bilinmeyen gibi düşünülür: çarpımı \(4\), toplamı \(-5\) \(\Rightarrow\) \(-1\) ve \(-4\).
\[x^4-5x^2+4=(x^2-1)(x^2-4)=(x-1)(x+1)(x-2)(x+2)\]
\end{cozum}

\begin{ornek}
\(x^3-9x=0\) \ ve \ \(3x^2+x-2=0\) denklemlerinin çözüm kümelerini bulunuz.
\end{ornek}
\begin{cozum}
\(x^3-9x=x(x-3)(x+3)=0\) \(\Rightarrow\) ÇK \(=\{-3,\ 0,\ 3\}\).\\
\(3x^2+x-2=(3x-2)(x+1)=0\) \(\Rightarrow\) \(x=\dfrac{2}{3}\) veya \(x=-1\): \ ÇK \(=\left\{-1,\ \dfrac{2}{3}\right\}\).
\end{cozum}

% =====================================================
\section{Küp Açılımları}

\begin{kural}[Özdeşlikler]
\[(a\pm b)^3=a^3\pm 3a^2b+3ab^2\pm b^3,\qquad a^3\pm b^3=(a\pm b)(a^2\mp ab+b^2)\]
\[a^3+b^3=(a+b)^3-3ab(a+b),\qquad a^3-b^3=(a-b)^3+3ab(a-b)\]
\end{kural}

\begin{dikkat}
\((a+b)^3\neq a^3+b^3\). \ \(a^2-ab+b^2\) ifadesi tam kare değildir; daha fazla ayrılmaz.
\end{dikkat}

\begin{ornek}
\((2x-1)^3\) açılımını yazınız.
\end{ornek}
\begin{cozum}
\((2x)^3-3(2x)^2\cdot 1+3(2x)\cdot 1^2-1^3=8x^3-12x^2+6x-1\).
\end{cozum}

\begin{ornek}
\(8x^3-27y^3\) ifadesini çarpanlarına ayırınız.
\end{ornek}
\begin{cozum}
\((2x)^3-(3y)^3=(2x-3y)\big((2x)^2+2x\cdot 3y+(3y)^2\big)=(2x-3y)(4x^2+6xy+9y^2)\).
\end{cozum}

\begin{ornek}
\(a+b=3\) ve \(a^3+b^3=9\) olduğuna göre \(a\cdot b\) kaçtır?
\end{ornek}
\begin{cozum}
\(a^3+b^3=(a+b)^3-3ab(a+b)\) \(\Rightarrow\) \(9=27-9ab\) \(\Rightarrow\) \(ab=2\).
\end{cozum}

\begin{ornek}
\(x-\dfrac{1}{x}=3\) olduğuna göre \(x^3-\dfrac{1}{x^3}\) kaçtır?
\end{ornek}
\begin{cozum}
\(a=x\), \(b=\frac{1}{x}\), \(ab=1\): \quad \(x^3-\dfrac{1}{x^3}=\Big(x-\dfrac{1}{x}\Big)^3+3\Big(x-\dfrac{1}{x}\Big)=27+9=36\).
\end{cozum}

\begin{ornek}
\(\dfrac{47^3+53^3}{100}\) işleminin sonucu kaçtır?
\end{ornek}
\begin{cozum}
\(47^3+53^3=(47+53)(47^2-47\cdot 53+53^2)\) ve \(47+53=100\):
\[\frac{47^3+53^3}{100}=2209-2491+2809=2527\]
\end{cozum}

% =====================================================
\section{Rasyonel İfadeleri Sadeleştirme}

\begin{kural}[Kural]
\(\dfrac{P}{Q}\) ifadesi yalnız \(Q\neq 0\) iken tanımlıdır. Pay ve payda çarpanlarına ayrılır, ortak \textbf{çarpanlar} sadeleştirilir. Tanım koşulu, sadeleştirmeden \textbf{önceki} paydalardan (bölmede ikinci kesrin payından da) yazılır.
\end{kural}

\begin{ornek}
\(\dfrac{x^2+5x+6}{x^2-9}\cdot\dfrac{x-3}{x+2}\) ifadesinin en sade hali nedir?
\end{ornek}
\begin{cozum}
\[\frac{(x+2)(x+3)}{(x-3)(x+3)}\cdot\frac{x-3}{x+2}=1\qquad(x\neq 3,\ x\neq -3,\ x\neq -2)\]
\end{cozum}

\begin{ornek}
\(\dfrac{2x^2-8}{x^2+4x+4}\) ifadesini sadeleştiriniz.
\end{ornek}
\begin{cozum}
\[\frac{2(x-2)(x+2)}{(x+2)^2}=\frac{2(x-2)}{x+2}=\frac{2x-4}{x+2}\qquad(x\neq -2)\]
\end{cozum}

\begin{ornek}
\(a=2027\) ve \(b=2026\) için \(\dfrac{a^3-b^3}{a^2+ab+b^2}+\dfrac{a^2-b^2}{a-b}\) ifadesinin değeri kaçtır?
\end{ornek}
\begin{cozum}
\(\dfrac{(a-b)(a^2+ab+b^2)}{a^2+ab+b^2}=a-b\), \quad \(\dfrac{(a-b)(a+b)}{a-b}=a+b\) \ (\(a\neq b\)).\\
Toplam: \((a-b)+(a+b)=2a=4054\).
\end{cozum}

% =====================================================
\section{Özet: Formül Kartı}
\begin{center}
\renewcommand{\arraystretch}{1.3}
\begin{tabular}{@{}ll@{}}
\toprule
Konu & Formül / Kural\\\midrule
Ortak çarpan & \(Ax+Ay=A(x+y)\); \ \(b-a=-(a-b)\)\\
Gruplandırma & \(ax+ay+bx+by=(x+y)(a+b)\)\\
İki kare farkı & \(a^2-b^2=(a-b)(a+b)\); \ \(a^2+b^2\) ayrılmaz\\
Tam kare & \((a\pm b)^2=a^2\pm 2ab+b^2\); \ \((a+b)^2-(a-b)^2=4ab\)\\
Kareler toplamı & \(a^2+b^2=(a+b)^2-2ab=(a-b)^2+2ab\)\\
\(x\) ve \(\frac{1}{x}\) & \(x^2+\frac{1}{x^2}=\big(x+\frac{1}{x}\big)^2-2\); \ \(x^3+\frac{1}{x^3}=\big(x+\frac{1}{x}\big)^3-3\big(x+\frac{1}{x}\big)\)\\
Eksik terim & \(x^2+bx+m\) tam kare \(\Leftrightarrow m=(b/2)^2\)\\
Üç terimli & \(x^2+bx+c=(x+m)(x+n)\), \ \(m+n=b\), \ \(mn=c\); \ \(ax^2+bx+c\): çapraz\\
Çarpım sıfır & \(A\cdot B=0\Rightarrow A=0\) veya \(B=0\)\\
Küp açılımı & \((a\pm b)^3=a^3\pm 3a^2b+3ab^2\pm b^3\)\\
Küp toplamı/farkı & \(a^3\pm b^3=(a\pm b)(a^2\mp ab+b^2)\)\\
Değer soruları & \(a^3+b^3=(a+b)^3-3ab(a+b)\); \ \(a^3-b^3=(a-b)^3+3ab(a-b)\)\\
Sadeleştirme & yalnız çarpanlar sadeleşir; tanım koşulu ilk paydadan\\\bottomrule
\end{tabular}
\end{center}

% =====================================================
\newpage
% Hafta 11 -- alıştırma gövdesi (ders notu ve ayrı alıştırma kağıdı tarafından \input edilir)
\section{Çözümlü Alıştırmalar}
Her soruyu önce kendiniz deneyin, sonra çözümü okuyun.

\begin{alistirma}
\(12a^2b-18ab^2\) ifadesini çarpanlarına ayırınız.
\end{alistirma}
\begin{cozum}
EBOB\((12,18)=6\); her iki terimde de \(a\) ve \(b\) var \(\Rightarrow\) ortak çarpan \(6ab\).
\[12a^2b-18ab^2=6ab\,(2a-3b)\]
Kontrol: \(6ab\cdot 2a=12a^2b\), \ \(6ab\cdot 3b=18ab^2\) \checkmark
\end{cozum}

\begin{alistirma}
\(x^2-y^2+3x-3y\) ifadesini çarpanlarına ayırınız.
\end{alistirma}
\begin{cozum}
İlk iki terim iki kare farkı, son iki terimde ortak çarpan \(3\):
\[x^2-y^2+3x-3y=(x-y)(x+y)+3(x-y)\]
Ortak parantez \((x-y)\) dışarı alınır:
\[=(x-y)(x+y+3)\]
\end{cozum}

\begin{alistirma}
\((x+1)^2-25\) ifadesini çarpanlarına ayırınız.
\end{alistirma}
\begin{cozum}
\(a=x+1\), \(b=5\) olmak üzere iki kare farkı:
\[(x+1)^2-5^2=(x+1-5)(x+1+5)=(x-4)(x+6)\]
Kontrol: \((x-4)(x+6)=x^2+2x-24\) ve \((x+1)^2-25=x^2+2x+1-25=x^2+2x-24\) \checkmark
\end{cozum}

\begin{alistirma}
\(a-b=5\) ve \(a^2+b^2=37\) olduğuna göre \(a\cdot b\) kaçtır?
\end{alistirma}
\begin{cozum}
\((a-b)^2=a^2-2ab+b^2\) özdeşliğinde verilenler yerine yazılır:
\[5^2=37-2ab\ \Rightarrow\ 25=37-2ab\ \Rightarrow\ 2ab=12\ \Rightarrow\ ab=6\]
(Kontrol: \(a=6,\ b=1\): \(a-b=5\), \(36+1=37\), \(ab=6\).)
\end{cozum}

\begin{alistirma}
\(x^2+kx+49\) bir tam kare ve \(k<0\) olduğuna göre \(k\) kaçtır?
\end{alistirma}
\begin{cozum}
\(x^2=(x)^2\), \(49=7^2\). Tam kare ifade \((x+7)^2\) veya \((x-7)^2\) olabilir.\\
\(k<0\) olduğundan \((x-7)^2=x^2-14x+49\) \(\Rightarrow\) \(k=-14\).
\end{cozum}

\begin{alistirma}
\(4x^2-4x-3\) ifadesini çarpanlarına ayırınız.
\end{alistirma}
\begin{cozum}
\begin{minipage}[t]{0.2\linewidth}
\begin{tabular}{ccc}
\(2x\) & & \(-3\)\\
\(2x\) & & \(1\)
\end{tabular}
\end{minipage}\hfill
\begin{minipage}[t]{0.62\linewidth}
\(4x^2=2x\cdot 2x\), \ \(-3=(-3)\cdot 1\).\\
Çapraz: \(2x\cdot 1+2x\cdot(-3)=2x-6x=-4x\) \checkmark
\[4x^2-4x-3=(2x-3)(2x+1)\]
\end{minipage}
\end{cozum}

\begin{alistirma}
\(a+b=6\) ve \(a\cdot b=5\) olduğuna göre \(a^3+b^3\) kaçtır?
\end{alistirma}
\begin{cozum}
\[a^3+b^3=(a+b)^3-3ab(a+b)=6^3-3\cdot 5\cdot 6=216-90=126\]
İkinci yol: \(a^2+b^2=36-10=26\); \ \(a^3+b^3=(a+b)(a^2-ab+b^2)=6\cdot(26-5)=126\).\\
(Kontrol: \(a=1,\ b=5\): \(1+125=126\).)
\end{cozum}

\begin{alistirma}
\(\dfrac{x^3+8}{x^2-4}\) ifadesini sadeleştiriniz ve tanım koşulunu yazınız.
\end{alistirma}
\begin{cozum}
Pay küp toplamı, payda iki kare farkı:
\[\frac{x^3+8}{x^2-4}=\frac{(x+2)(x^2-2x+4)}{(x-2)(x+2)}=\frac{x^2-2x+4}{x-2}\]
Tanım koşulu ilk paydadan: \(x^2-4\neq 0\) \(\Rightarrow\) \(x\neq 2\) ve \(x\neq -2\).
\end{cozum}

\kendinizoncesi
\section{Alıştırmalar (Kendiniz Çözün)}

\begin{enumerate}[label=\textbf{\arabic*.},leftmargin=*,itemsep=6pt]
\item \(7x^2-14x\) ifadesini çarpanlarına ayırınız.
\item \(xy+2x+3y+6\) ifadesini çarpanlarına ayırınız.
\item \(4x^2-81\) ifadesini çarpanlarına ayırınız.
\item \(75^2-25^2\) işleminin sonucu kaçtır?
\item \(x^4-81\) ifadesini çarpanlarına ayırınız.
\item \((2a+5)^2\) açılımını yazınız.
\item \(a+b=9\) ve \(a\cdot b=14\) olduğuna göre \(a^2+b^2\) kaçtır?
\item \(x+\dfrac{1}{x}=6\) olduğuna göre \(x^2+\dfrac{1}{x^2}\) kaçtır?
\item \(x^2-10x+m\) bir tam kare ise \(m\) kaçtır?
\item \(x^2+5x-14\) ifadesini çarpanlarına ayırınız.
\item \(2x^2-5x-3\) ifadesini çarpanlarına ayırınız.
\item \(x^2-6x+8=0\) denkleminin çözüm kümesini bulunuz.
\item \(x^3-27\) ifadesini çarpanlarına ayırınız.
\item \(a-b=1\) ve \(a\cdot b=6\) olduğuna göre \(a^3-b^3\) kaçtır?
\item \(\dfrac{x^2+2x-8}{x^2-4}\) ifadesini sadeleştiriniz ve tanım koşulunu yazınız.
\item \(x+\dfrac{1}{x}=4\) olduğuna göre \(x^3+\dfrac{1}{x^3}\) kaçtır?
\end{enumerate}

\needspace{10\baselineskip}
\subsection*{Cevap Anahtarı}
{\small
\begin{multicols}{2}
\begin{enumerate}[label=\textbf{\arabic*.},leftmargin=*,itemsep=1pt]
\item Ortak çarpan \(7x\): \(7x(x-2)\).
\item \(x(y+2)+3(y+2)=(y+2)(x+3)\).
\item \((2x)^2-9^2=(2x-9)(2x+9)\).
\item \((75-25)(75+25)=50\cdot 100=5000\).
\item \((x^2-9)(x^2+9)=(x-3)(x+3)(x^2+9)\).
\item \(4a^2+20a+25\).
\item \(9^2-2\cdot 14=81-28=53\).
\item \(6^2-2=34\).
\item \((x-5)^2\Rightarrow m=25\).
\item Çarpım \(-14\), toplam \(5\): \((x+7)(x-2)\).
\item Çapraz: \((2x+1)(x-3)\).
\item \((x-2)(x-4)=0\Rightarrow\{2,\,4\}\).
\item \((x-3)(x^2+3x+9)\).
\item \(1^3+3\cdot 6\cdot 1=19\).
\item \(\dfrac{(x+4)(x-2)}{(x-2)(x+2)}=\dfrac{x+4}{x+2}\); \ \(x\neq\pm 2\).
\item \(4^3-3\cdot 4=64-12=52\).
\end{enumerate}
\end{multicols}}


\vspace{6pt}
\noindent\textcolor{anamavi}{\rule{\textwidth}{0.5pt}}\\
{\small Gelecek hafta: Oran ve orantı, ortalamalar.}

\end{document}
